% ============================================================
\section{Introduction}
\label{sec:introduction}
% ============================================================

Rule-based tactical asset allocation (TAA) strategies are popular precisely
because they are transparent: an allocation is driven by a handful of explicit
rules (a momentum signal, a ranking cutoff, an absolute-trend filter, a
portfolio-construction criterion) rather than by an opaque optimiser.
Momentum, the workhorse signal of these rules, is one of the most robust
regularities in asset markets, both cross-sectionally and as a time-series effect
across asset classes \citep{jegadeesh1993, moskowitz2012, asness2013}.
Yet that same transparency raises a question the literature rarely answers head on.
Among the several rules a typical strategy stacks, which ones actually generate the
out-of-sample edge, which are inert, and is the whole thing anything more than a
repackaged time-series-momentum (trend) premium? Without an answer, a favourable
backtest is hard to tell apart from the product of a parameter search
\citep{yang2019, harvey2018}.

This paper takes one concrete cross-asset momentum strategy and treats it as a case
study for exactly that question. We derive each rule on a locked training window and
then decompose its out-of-sample performance rule by rule. The tools are a
matched-permutation (skill-versus-luck) test and a nested spanning-alpha ladder, and
we then put the residual edge through controls for known risk factors (including a
trend factor built from the same assets), transaction costs, a multiple-testing
adjustment, and a replication on 100\% real data. The strategy itself is designed
under a hard drawdown-control objective, which we keep as the evaluation criterion.
But the contribution is the attribution, not the strategy.

The strategy under study combines four components.
First, a multi-period (13612U) momentum score averages 1-, 3-, 6-, and 12-month
signals to reduce dependence on any single lookback horizon.
Second, a top-$7$ breadth pre-filter restricts the candidate set to the
higher-momentum half of the universe.
Third, an absolute-momentum threshold defines an eligible set and drives a
breadth-based defensive switch: when fewer than three assets clear the threshold,
the cross-section of momentum has broadly deteriorated and the portfolio rotates
to long- and intermediate-duration treasuries, or to cash if bond momentum is also
negative \citep{keller2017, antonacci2014}.
Fourth, a minimum-correlation triplet selection rule implements the central
diversification insight of modern portfolio theory \citep{markowitz1952} by
holding the three lowest-correlation assets among the eligible candidates.
Regime detection is thus driven by the \emph{breadth of momentum} (the count of
assets clearing the absolute threshold) rather than by a separate market-return
trigger.

We evaluate the strategy on a universe of 15 diversified ETFs covering global
equities, real estate, commodities, and fixed income; exchange-traded funds make
systematic multi-asset rotation operationally feasible by providing liquid,
low-cost exposure to broad asset classes.
Because the ETFs were launched between 2002 and 2008, a real-ETF backtest is too
short for a powered train/validation/test design; we therefore extend each series
back to 2003 using documented proxies (underlying indices or OLS-calibrated
correlated instruments), with the data-construction choices and their limitations
reported in Section~\ref{sec:methodology}.
The empirical design uses a strict nested temporal split: the momentum signal, the
breadth count, and the triplet size are derived on the 2004--2015 training window
(each lands on its ex-ante literature value) and confirmed on a 2016--2020
validation window, while the absolute threshold is fixed ex ante at the risk-free
rate rather than optimised; the frozen strategy is then evaluated once on a locked
2021--2025 test block. Supporting robustness analyses are reported over the broader
post-training sample (2016--2025) for statistical power.
The out-of-sample windows span the COVID-19 crash (validation) and the 2022
inflation and rate shock and 2023--2025 recovery (test).

The main result is that the strategy delivers competitive risk-adjusted returns
with materially lower drawdowns than passive benchmarks.
Over the full 2004--2025 period it achieves a CAGR of 11.00\%, a Sharpe ratio of
0.892, and a maximum drawdown of $-11.86\%$.
On the locked test block (2021--2025) it achieves a CAGR of 11.42\%, a Sharpe ratio
of 0.873 (the highest of the benchmark panel), and a maximum drawdown of $-7.06\%$,
roughly one-third that of $1/N$ ($-23.8\%$), 60/40 ($-21.1\%$), and the S\&P~500
($-24.8\%$).

The paper makes three contributions.
The central one is a transparent, two-pronged attribution
framework for rule-based TAA strategies, deployed here on a concrete four-rule
strategy as a case study. The framework combines a matched-permutation
(skill-versus-luck) test, which replaces each rule's informed decision with a random
draw from the same feasible set, with a nested incremental spanning-alpha ladder
that isolates each rule's marginal out-of-sample contribution. The resulting
decomposition is more informative, and more defensible against
data-snooping charges, than either a single in-sample optimisation or an aggregate
out-of-sample test. Applied to GAT, it finds that the momentum rule beats its
matched random baseline ($p = 0.0025$); the minimum-correlation rule contributes
through variance reduction (variance ratio $1.28$, $p = 0.010$) rather than
through return; and the remaining parameters, fixed ex ante in flat regions of the
objective, are statistically indistinguishable from random matched choices. The
incremental ladder tells the same story. No single rule is individually significant
in isolation; the edge emerges from stacking complementary rules, and the
combination is what generates the alpha.
The second contribution is a formal and empirical account of the minimum-correlation
construction. Proposition~\ref{prop:mincorr} shows analytically that, within a
momentum-screened eligible set with homogeneous expected returns, the
minimum-average-correlation triplet is the minimum-variance equal-weight portfolio.
The data then show that this is precisely its empirical role. It compresses
portfolio variance and tail risk at statistically equal expected
return and without inverting a covariance matrix, so the gain is risk control rather
than return. Reporting a null return contribution
alongside a significant risk contribution is, we argue, an honest finding with direct
implications for anyone designing multi-rule TAA systems: the rule earns its
place through drawdown compression, not alpha.
The third contribution is breadth of evidence for the spanning alpha.
The locked test block delivers $+6.9\%$ per annum against $1/N$ ($p = 0.008$) and
a residual of $+4$--$5\%$/yr against joint controls for equity, duration, and a
long-only time-series-momentum factor built from the same 15 assets. That long-only
factor is the hardest appropriate control for a long-only strategy, and GAT still
clears it ($+7.6\%$/yr against
TSMOM$_{\text{LO}}$, $p = 0.018$; $+4.3\%$ against the standard long-short
TSMOM, $p = 0.04$, both on the locked test). The rest of the robustness battery
points the same way. Realistic transaction costs (break-even
$\approx 62$ bps against $10\times$/yr turnover), temporal stability, universe
perturbation, a Deflated Sharpe ratio of $0.975$, and a replication on 100\% real
ETF data each fail to overturn the result. We also confront the strategy's main
practical weakness, its tax inefficiency relative to buy-and-hold: the tax drag is
real ($\approx 3$ percentage points of annual return), yet GAT retains the highest
after-tax Sharpe of the benchmark panel (Section~\ref{sec:res_tax}).
The strategy itself is transparent, reproducible, and implementable with monthly
ETF trades, but our central message is methodological: the attribution framework
reveals where a rule-based TAA edge comes from and, equally important, where it
does not.

The remainder of the paper is structured as follows.
Section~\ref{sec:literature} reviews the relevant literature.
Section~\ref{sec:methodology} describes the strategy design and data.
Section~\ref{sec:results} reports in-sample performance, the rule-by-rule
attribution, out-of-sample and full-period results, and the robustness battery.
Section~\ref{sec:discussion} discusses what the attribution implies about the
source of the edge.
Section~\ref{sec:conclusion} concludes.
